\documentclass[10pt]{article}
\usepackage[margin=0.75in]{geometry}
\usepackage{xcolor}
\usepackage{tcolorbox}
\usepackage{hyperref}
\usepackage{enumitem}
\usepackage{booktabs}

% Define colors
\definecolor{osaorange}{RGB}{220,115,38}
\definecolor{aiblue}{RGB}{30,144,255}

% Configure hyperref
\hypersetup{
    colorlinks=true,
    linkcolor=blue,
    urlcolor=blue
}

% Compact spacing
\setlength{\parindent}{0pt}
\setlength{\parskip}{4pt}
\setlist{noitemsep,topsep=2pt}

\begin{document}

% Orange header
\begin{tcolorbox}[colback=osaorange,colframe=black,boxrule=2pt,arc=0pt,left=6pt,right=6pt,top=8pt,bottom=8pt]
\centering
\textcolor{white}{\textbf{\large WEEK 9 ASSIGNMENT}}\\
\textcolor{white}{\textbf{H524 Introduction to Biostatistics - Fall 2025}}
\end{tcolorbox}

\vspace{4pt}

\noindent\textbf{Total Points:} 10 \quad|\quad \textbf{Due:} Sunday 11:59 PM

\vspace{6pt}

\section*{Part A: Multiple Choice Questions (7 points)}

Answer the following 7 questions on correlation and simple linear regression. Each question is worth 1 point.

\subsection*{Question 1: Correlation Coefficient Range (1 point)}
The Pearson correlation coefficient (r) ranges from:
\begin{enumerate}[label=\Alph*., leftmargin=20pt]
\item 0 to 1
\item -1 to 0
\item -1 to 1
\item -$\infty$ to +$\infty$
\end{enumerate}

\subsection*{Question 2: Interpreting Correlation (1 point)}
A correlation coefficient of r = -0.85 indicates:
\begin{enumerate}[label=\Alph*., leftmargin=20pt]
\item A weak positive linear relationship
\item A strong positive linear relationship
\item A weak negative linear relationship
\item A strong negative linear relationship
\end{enumerate}

\subsection*{Question 3: Correlation vs. Causation (1 point)}
If two variables have a strong positive correlation (r = 0.90), this means:
\begin{enumerate}[label=\Alph*., leftmargin=20pt]
\item One variable causes changes in the other
\item The variables tend to increase together, but this does not prove causation
\item There is no relationship between the variables
\item The regression line has a slope of 0.90
\end{enumerate}

\subsection*{Question 4: Simple Linear Regression Equation (1 point)}
In the regression equation $\hat{y} = a + bx$, what does ``b'' represent?
\begin{enumerate}[label=\Alph*., leftmargin=20pt]
\item The y-intercept
\item The slope of the regression line
\item The correlation coefficient
\item The predicted value of y
\end{enumerate}

\subsection*{Question 5: Coefficient of Determination (1 point)}
If the correlation coefficient r = 0.60, what is the coefficient of determination (R²)?
\begin{enumerate}[label=\Alph*., leftmargin=20pt]
\item 0.36
\item 0.60
\item 0.77
\item 1.67
\end{enumerate}

\subsection*{Question 6: Interpreting R² (1 point)}
An R² value of 0.64 means:
\begin{enumerate}[label=\Alph*., leftmargin=20pt]
\item 64\% of the variance in y is explained by x
\item 64\% of the variance in x is explained by y
\item The correlation is r = 0.64
\item The slope of the regression line is 0.64
\end{enumerate}

\subsection*{Question 7: Regression Assumptions (1 point)}
Which of the following is NOT an assumption of simple linear regression?
\begin{enumerate}[label=\Alph*., leftmargin=20pt]
\item Linearity (relationship between x and y is linear)
\item Independence of observations
\item Normality of residuals
\item The independent variable (x) must be normally distributed
\end{enumerate}

\newpage

\section*{Part B: Group Project Check-In - Presentation Readiness (3 points)}

\begin{tcolorbox}[colback=aiblue!10,colframe=aiblue,boxrule=1pt,arc=2pt,left=6pt,right=6pt,top=6pt,bottom=6pt]
\textbf{GROUP PROJECT: FINAL CHECK BEFORE PRESENTATIONS}
\end{tcolorbox}

\vspace{4pt}

\textbf{Reminder:} Presentations are NEXT WEEK (Week 10)!
\begin{itemize}[leftmargin=20pt]
\item \textbf{Monday 12-1:20pm:} Groups 1 \& 2
\item \textbf{Wednesday 12-1:20pm:} Groups 3 \& 4
\end{itemize}

\subsection*{Your Assignment}

\textbf{Group Leader submits} evidence that your presentation is ready:

\textbf{Choose ONE of the following to submit:}

\textbf{Option A: Presentation Outline} (preferred)
\begin{itemize}[leftmargin=20pt]
\item List of slides/sections (bullet points)
\item Who will present each section
\item Timing estimate for each part
\end{itemize}

\textbf{Option B: Draft Slide Deck}
\begin{itemize}[leftmargin=20pt]
\item PowerPoint, Google Slides, or PDF
\item Doesn't need to be final, but should show structure
\end{itemize}

\textbf{Option C: Presentation Plan}
\begin{itemize}[leftmargin=20pt]
\item Written description of what you'll cover
\item Division of speaking roles
\item Confirmation that all members will participate
\end{itemize}

\subsection*{Submission Format}

\textbf{Submit as:} File upload (PDF, PPT, DOCX) OR text entry

\textbf{Length:}
\begin{itemize}[leftmargin=20pt]
\item Outline: 1 page
\item Slides: Any number (draft is fine)
\item Plan: 1-2 paragraphs
\end{itemize}

\textbf{Grading:} 3 points for demonstrating presentation preparation

\subsection*{Presentation Reminders}

\textbf{Your presentation should include:}
\begin{itemize}[leftmargin=20pt]
\item Brief dataset overview
\item Key statistical findings
\item \textbf{AI comparison highlights} (30-40\% of presentation!)
\item Conclusions and what you learned
\end{itemize}

\textbf{Requirements:}
\begin{itemize}[leftmargin=20pt]
\item 15 minutes (+ 5 min Q\&A)
\item All team members participate
\item Use slides
\item Attend your assigned day (Mon or Wed)
\end{itemize}

\textbf{Time to finalize:}
\begin{itemize}[leftmargin=20pt]
\item This week: Draft slides and practice
\item Next week: Final touches and delivery
\end{itemize}

\subsection*{Notes}

\begin{itemize}[leftmargin=20pt]
\item \textbf{Only the group leader submits}
\item This check-in ensures you're ready for Week 10
\item If you submit an outline now, you can refine slides this week
\item Don't wait until the day before to start your presentation!
\item \textbf{Grading (10\% of course grade) is next week} - take this seriously!
\end{itemize}

\end{document}
